% ch9_d 第9章 §9.2尾–§9.3 (书页 393–405 / p408–p420)
% 块边界判定：
%   起点——书页 393（p408）顶部为图 9.5（属 9.2.7 节 T 变换讨论，前块止于书页 392 末的
%     式 (9.2.44)），页面首段 "The above describes ..." 为完整新句，无前块溢出内容，
%     故本块不设 %JOIN%，亦无需 %--C-TAIL-- 区。
%   终点——书页 405（p420）末句 "The key to obtaining the most general Z2A0013 ansatz
%     is to" 延续至书页 406（p421）首行 "find out which $u_m^\mu$ must vanish ..."。
%     按约定，半句译文（"……的关键是"）留在本文件正文最后，不写入任何 TAIL 区，
%     由 ch9_e 在其正文开头以 %JOIN% 续接。
% 蓝本问题记录：式 (9.4.11) 蓝本漏排 $G_y(\pmb{i})=\tau^0$ 一栏且上标误作 i_y（以图为准）；
%   式 (9.4.13) 蓝本作 $(-)^{m_y i_x}$，图上为 $(-)^{m i_x}$；书页 402 "±G_yT_x" 为
%   "±G_yT_y" 之误（据脚注 70 与式 (9.4.6)），已按本意译出并加译注。

%FIG{9.5}: {跨越 $x$ 线与 $y$ 线的链获得一个额外的负号。}

上述结果描述了平均场拟设在 $T$ 变换下的变换方式。自旋哈密顿量 (9.1.1) 与平均场哈密顿量 (9.2.4) 在时间反演变换下都是不变的。

\subsection*{问题 9.2.8}

试求在真实的时间反演变换 $S\to-S$ 下，$U_{ij}$ 与 $a_0^l(\pmb{i})$ 如何变换。

\section{有能隙自旋液体态中的拓扑序}

\begin{keypoints}
\item 自旋液体基态的拓扑简并。
\item 稳健的拓扑简并基态的存在意味着拓扑序的存在。
\end{keypoints}

我们用平均场理论构造了 $Z_2$ 有能隙态与手征自旋态。这两个自旋液体态的所有激发都有有限的能隙，并且保持平移对称性。我们想指出，这里研究的手征自旋液体与 $Z_2$ 有能隙自旋液体包含一些新型的序，它们既不能用对称性、也不能用与破缺对称性相联系的序参量来刻画。对 $Z_2$ 有能隙态而言，这一点相当明显，因为它不破缺任何对称性。手征自旋液体自发破缺了时间反演对称性与宇称对称性，并有相应的局域序参量来描述这种对称性破缺。正如我们后面将看到的，手征自旋液体还包含一些无法用对称性破缺描述的额外结构。由于手征自旋液体与 $Z_2$ 有能隙自旋液体都有有限的能隙，这类新型的序就是拓扑序。

正如 8.2.1 节所指出的，证明 $Z_2$ 有能隙态与手征自旋态具有非平庸拓扑序的一种办法，是证明这两个态具有拓扑简并的基态。让我们先考虑 $Z_2$ 有能隙态及其基态简并度。

我们把 $Z_2$ 有能隙态放到一个环面上，$x$ 与 $y$ 两个方向的格点数都取偶数。考虑由一个平均场解 $\bar{u}_{\pmb{ij}}$ 构造出的下列四个拟设：
\begin{equation*}
u_{\pmb{ij}}^{(m,n)}=(-)^{m s_x(\pmb{ij})}(-)^{n s_y(\pmb{ij})}\bar{u}_{\pmb{ij}}
\tag{9.3.1}
\end{equation*}
其中 $m,n=0,1$。这里 $s_{x,y}(\pmb{ij})$ 取 0 或 1：若链 $\pmb{ij}$ 跨越 $x$ 线，则 $s_x(\pmb{ij})=1$（见图 9.5），否则 $s_x(\pmb{ij})=0$；类似地，若链 $\pmb{ij}$ 跨越 $y$ 线，则 $s_y(\pmb{ij})=1$，否则 $s_y(\pmb{ij})=0$。

我们注意到，不同 $m$、$n$ 的 $u_{\pmb{ij}}^{(m,n)}$ 是局域规范等价的。这是因为，在无限系统上，改变 $u_{\pmb{ij}}\to(-)^{m s_x(\pmb{ij})}(-)^{n s_y(\pmb{ij})}u_{\pmb{ij}}$ 可以由一个 $SU(2)$ 规范变换 $u_{\pmb{ij}}\to W_{\pmb{i}}u_{\pmb{ij}}W_{\pmb{j}}^{\dagger}$ 生成，其中 $W_{\pmb{i}}=(-)^{m\Theta(i_x)}(-)^{n\Theta(i_y)}$，而 $\Theta(n)=1$ 若 $n>0$，$\Theta(n)=0$ 若 $n\leqslant0$。结果，不同的拟设 $u_{\pmb{ij}}^{(m,n)}$ 穿过每个方格的 $SU(2)$ 通量都相同。

对大系统而言，能量是 $u_{\pmb{ij}}$ 的局域函数，并且若 $u_{\pmb{ij}}$ 与 $\tilde{u}_{\pmb{ij}}$ 规范等价，则有 $F(u_{\pmb{ij}})=F(\tilde{u}_{\pmb{ij}})$。因此，不同 $u_{\pmb{ij}}^{(m,n)}$ 的能量相同（在热力学极限下）。另一方面，不同 $m$、$n$ 的 $u_{\pmb{ij}}^{(m,n)}$ 在整体意义上并不规范等价。一个自旋子沿 $x$（或 $y$）方向绕环面一整圈，会获得相位 $\mathrm{e}^{\mathrm{i}m\pi}$（或 $\mathrm{e}^{\mathrm{i}n\pi}$）。因此，$u_{\pmb{ij}}^{(mn)}\big|_{m,n=0,1}$ 描述不同的、相互正交的简并基态（另见问题 9.3.1）。这样我们就发现，$Z_2$ 有能隙态在环面上有四个简并基态（Read and Chakraborty, 1989; Wen, 1991a）。从自旋子绕环面一圈所获得的相位可以看出，$m$ 与 $n$ 标记穿过环面两个孔的 $\pi$ 通量的数目（见图 9.6）。四个简并基态对应于把 $\pi$ 通量穿入两个孔的不同方式。这一图像可以推广到更高亏格的黎曼面。简并基态可以这样构造：让 $\pi$ 通量以零个或一个单位穿过亏格为 $g$ 的黎曼面上的 $2g$ 个孔（见图 9.6）。这给出亏格为 $g$ 的黎曼面上 $2^{2g}$ 个简并基态。这种在不同黎曼面上的简并图案是 $Z_2$ 规范理论的标志。这些简并度与 $Z_2$ 涡旋激发的存在，以一种物理的方式表明，$Z_2$ 有能隙态的低能有效理论确实是一个 $Z_2$ 规范理论（见问题 9.3.1）。把自旋子包括进来之后，$Z_2$ 有能隙态的一阶平均场理论描述的是耦合到 $Z_2$ 规范理论上的费米子。$Z_2$ 有能隙态的物理性质可以从这样一个有效理论得到。

%FIG{9.6}: {(a) 亏格 $g=1$ 的黎曼面（环面）上有 $2g=2$ 个孔；(b) 亏格 $g=2$ 的黎曼面上有 $2g=4$ 个孔。我们可以往这些孔里穿入零通量或 $\pi$ 通量。这给出 $Z_2$ 有能隙自旋液体的 $2^{2g}=4$ 或 16 个简并基态。}

在一个有限的紧致化格点上，系统从一个基态隧穿到另一个基态的唯一途径是如下隧穿过程：先产生一对 $Z_2$ 涡旋，其中一个 $Z_2$ 涡旋绕环面一整圈，然后与另一个 $Z_2$ 涡旋湮没。这样的过程实质上给环面的孔添加了一个单位的 $\pi$ 通量，并使 $m$ 或 $n$ 改变 1。由于通量守恒，不同的基态无法通过任何局域涨落相互隧穿。这一结果的直接推论是，有限格点上不同基态之间的能量劈裂预期为 $\mathrm{e}^{-L/\xi}$ 的量级，其中 $1/\xi$ 与 $Z_2$ 涡旋的有限能隙相关。

在平均场理论中，基态的简并是规范不变性的一个后果。即使我们在原来的自旋哈密顿量中加入下列任意微扰，规范不变性仍然严格成立：
\begin{equation*}
\delta H=\sum\delta J_{ij}\,\pmb{S}_i\cdot\pmb{S}_j+\cdots
\tag{9.3.2}
\end{equation*}
其中 $\delta H$ 可以破坏平移对称性、旋转对称性等。因此上述论证依然成立，即使对修改后的哈密顿量，平均场基态仍保持四重简并。我们预期这一结果在平均场理论之外也成立。只要微扰足够弱、不至于关闭 $Z_2$ 涡旋的能隙，任何微扰都无法改变 $Z_2$ 有能隙态的基态简并度。因此，基态简并度是刻画自旋液体态不同相的一个普适量子数。这也表明 $Z_2$ 有能隙态中存在非平庸的拓扑序。

现在让我们考虑手征自旋态中的拓扑序。同样，我们想计算手征自旋态在比如环面上的基态简并度。计算基态简并度最容易的办法，是使用有效 Chern--Simons 理论 (9.1.28)，它描述 $a_{\mu}$ 集体模式的动力学。$U(1)$ Chern--Simons 理论的基态简并度已在 FQH 态的情形中讨论过。重复 8.2.1 节的计算，我们发现基态简并度是 2。我们想指出，这一二重简并是针对破缺 $T$ 与 $P$ 的基态的某一个扇区（sector）而言的，例如 $\langle\pmb{S}_1\cdot(\pmb{S}_2\times\pmb{S}_3)\rangle>0$ 的扇区。另一个扇区 $\langle\pmb{S}_1\cdot(\pmb{S}_2\times\pmb{S}_3)\rangle<0$ 也有两个简并基态。因此，总的基态简并度是 $4=2\times2$：一个因子 2 来自 $T$、$P$ 破缺，另一个因子 2 来自规范涨落。

由规范涨落引起的二重基态简并，源于手征自旋态中的非平庸拓扑序。这一简并依赖于空间的拓扑。在亏格为 $g$ 的黎曼面上，规范涨落给出 $2^g$ 个简并基态。这一简并度同样直接与规范结构相关，并且可以证明它对自旋哈密顿量的任意局域微扰都是稳健的。正如 FQH 态一样，手征自旋态的基态简并度与自旋子的分数统计密切相关。一般地可以证明，若自旋子具有分数统计 $\theta=\frac{p\pi}{q}$，则手征自旋态在亏格为 $g$ 的黎曼面上将有 $q^g$ 重简并的基态（Wen, 1990）。此外，作为非平庸拓扑序的另一个结果，手征自旋态具有无能隙的手征边缘激发（Wen, 1992）。

为了看出拓扑序概念的用处，让我们考虑下列物理问题：手征自旋态与 $Z_2$ 有能隙态之间的差别是什么？有人可能马上会说，这两个态有不同的对称性。然而，如果我们修改哈密顿量使 $T$ 与 $P$ 对称性破缺，两个态就会有相同的对称性。这时我们仍然可以问这两个态是否相同——即一个态能否不经相变而连续变形为另一个态。当 $T$ 与 $P$ 对称性被显式破缺时，手征自旋态有两个简并基态，而 $Z_2$ 有能隙态在环面上仍有四个简并基态。因此，即使具有相同的对称性，手征自旋态与 $Z_2$ 有能隙态也是不同的。两态的差别在于它们有不同的基态简并度和不同的拓扑序。

\subsection*{问题 9.3.1}

\textbf{非共线 $SU(2)$ 通量态中的 $Z_2$ 规范结构}

设 $u_{\pmb{ij}}=s_{ij}\bar{u}_{\pmb{ij}}$，其中 $s_{ij}=\pm1$，$\bar{u}_{\pmb{ij}}$ 是一个具有非共线 $SU(2)$ 通量的拟设。令 $|\{s_{ij}\}\rangle$ 为 $|\Psi_{\text{mean}}^{(u_{\pmb{ij}})}\rangle$ 的投影（见式 (9.2.11)），即
\begin{equation*}
|\{s_{ij}\}\rangle=\langle0_{\psi}|\prod_n\psi_{1,\pmb{i}_n}\psi_{2,\pmb{i}_n}|\Psi_{\text{mean}}^{(u_{\pmb{ij}})}\rangle.
\end{equation*}
\begin{enumerate}[label=(\alph*)]
\item 证明：与 $Z_2$ 规范理论中的态一样，若 $s_{ij}$ 与 $\tilde{s}_{ij}$ 通过一个 $Z_2$ 规范变换相联系，则 $|\{s_{ij}\}\rangle=|\{\tilde{s}_{ij}\}\rangle$。
\item 证明：式 (9.3.1) 所定义的四个 $u_{\pmb{ij}}^{(m,n)}$（其中 $\bar{u}_{\pmb{ij}}$ 为具有非共线 $SU(2)$ 通量的一般拟设）彼此不 $SU(2)$ 等价。（提示：考虑通过三个具有同一基点的回路 $C_{1,2,3}$ 的 $SU(2)$ 通量 $P(C_1)$、$P(C_2)$ 与 $P(C_3)$。这里 $C_1$ 与 $C_2$ 是两个小回路，满足 $[P(C_1),P(C_2)]\neq0$；$C_3$ 是绕环面一圈的大回路。）
\end{enumerate}

\section{对称自旋液体中的量子序}

\begin{keypoints}
\item 量子序是量子相中的非对称性破缺序。
\item 量子序既适用于有能隙量子态，也适用于无能隙量子态。
\end{keypoints}

在这一节中，我们将基于平均场拟设的不变群发展一种量子序理论。我们将把这个不变群称为投影对称群（projective symmetry group，PSG）。我们的理论使我们能够刻画并对数百种具有相同对称性的自旋液体态进行分类。与拓扑序的基态简并度不同，投影对称群既能刻画有能隙自旋液体，也能刻画无能隙自旋液体。

我们已经看到，可以有许多具有相同对称性的不同自旋液体。其中许多自旋液体占据相空间中的一个有限区域，代表稳定的量子相。于是我们在这里面临的处境与量子霍尔效应中的情形类似：存在许多无法用对称性和序参量区分的量子相。量子霍尔液体有有限的能隙，我们可以用基态简并度或更一般的拓扑序来描述这些不同态的内部序。这里我们仍然可以用拓扑序来描述有能隙自旋液体的内部序。然而，我们还有许多其他稳定的量子自旋液体，它们具有无能隙激发。

为了描述无能隙量子自旋液体（以及有能隙自旋液体）的内部序，我们在第 8 章中引入了一个新概念——量子序（quantum order）——它描述量子相中的非对称性破缺序（见图 8.7）。引入量子序的关键点是：量子相一般不能完全用破缺的对称性和局域序参量来刻画。量子霍尔态以及前几节构造的稳定自旋液体都说明了这一点。然而，要使量子序的概念有用，我们需要找到量子序的具体数学刻画。由于量子序既不由对称性也不由序参量描述，我们需要找到一种全新的刻画方式。Wen (2002c) 提出用投影对称群来刻画量子自旋液体中的量子（或拓扑）序。投影对称群的提出受到 9.2.6 节如下观察的启发：不同对称自旋液体的拟设，在同样的对称性变换之后跟上\emph{不同}的规范变换下不变。我们可以利用这些不同的规范变换来区分具有相同对称性的不同自旋液体。在接下来两节中，我们将在一般而形式化的框架下引入投影对称群。

\subsection{量子序与普适性质}

\begin{keypoints}
\item 定义一种新型的序，就是要找到一种新的普适性质。
\item （某一类）量子序可以用投影对称群描述。
\item 投影对称群是一个平均场拟设的不变群（或“对称性”群）。
\end{keypoints}

我们知道，定义一个量子相，就是要找出基态波函数的普适性质。我们可以利用普适性质把基态波函数归入若干普适类，使得每一类中的波函数具有相同的普适性质。这些普适类将对应于量子相。

然而，要找出多体波函数的普适性质是非常困难的。让我们考虑自由费米系统中的量子序（或普适类），以对这种困难获得一些直观认识。我们知道，自由费米子基态由一个 $N$ 个变量的反对称波函数描述。该反对称函数具有 Slater 行列式的形式，即 $\Psi(x_1,\ldots,x_N)=\det(M)$，其中 $M$ 的矩阵元由 $M_{mn}=\psi_n(x_m)$ 给出，而 $\psi_n$ 是单费米子波函数。寻找自由费米系统中量子序的第一步，是找到一种合理的方式把 Slater 行列式波函数归类。如果我们只知道实空间的多体函数 $\Psi(x_1,\ldots,x_N)$，这将是非常困难的。然而，如果我们用傅里叶变换把实空间波函数变到动量空间的波函数，就可以按照费米面拓扑把不同的波函数归类。要真正证明上述分类是合理的并对应于不同的量子相，我们需要证明：当基态波函数从一个类变到另一个类时，基态能量在转变点处总会出现奇异性。这就引出了我们对自由费米系统中量子序的理解（见 8.3.2 节）。费米面拓扑是一种普适性质，它使我们能够对不同的费米子波函数进行分类，并刻画自由费米子的不同量子相。这里我们想强调，如果没有傅里叶变换，就很难从实空间多体函数 $\Psi(x_1,\ldots,x_N)$ 中看出费米面拓扑。

为了理解自旋液体中的量子序，我们需要找到一种办法把自旋液体波函数 $\Psi_{\text{spin}}(\pmb{i}_1,\ldots,\pmb{i}_{N_t})$ 归类，其中 $\pmb{i}_m$ 标记自旋朝上的位置。这里缺少的是相应的“傅里叶”变换。正如费米面的拓扑一样，很难直接从实空间波函数中看出普适性质（如果有的话）。目前，我们还不知道如何对自旋液体波函数 $\Psi_{\text{spin}}$ 进行分类。

不过，我们不想就此放弃。受投影构造的启发，这里我们考虑一个较简单的问题来入手。我们把自己限制在可以由拟设 $(u_{\pmb{ij}},a_0^l\tau^l)$ 经式 (9.2.11) 得到的那一类多体波函数之内。我们不去寻找一般多体波函数的普适性质，而是设法找出这一子类中多体波函数的普适性质。由于这一子类中的多体波函数由拟设 $(u_{\pmb{ij}},a_0^l\tau^l)$ 标记，波函数的普适性质实际上就是拟设的普适性质，这大大简化了问题。诚然，有人可能会反驳说，拟设（或该子类波函数）的普适性质未必是自旋量子相的普适性质。对某些拟设来说情况确实如此。然而，如果拟设 $(u_{\pmb{ij}},a_0^l\tau^l)$ 所描述的平均场态对任何涨落都是稳定的（即在平均场态附近的任何涨落都不会引起红外发散），那么该平均场态就忠实地描述一个自旋量子态，拟设的普适性质也就是相应自旋量子相的普适性质。这样便在拟设的性质与物理自旋液体的性质之间建立起联系。

那么，拟设的普适性质应当是什么呢？受朗道对称性破缺序理论的启发，我们想在此提出：拟设的不变群（或“对称性”群）是拟设的一种普适性质。这个群将被称为投影对称群（PSG）。我们将在后面论证，某些 PSG 确实是量子相的普适性质，并且这些 PSG 可以用来刻画量子自旋液体中的量子序。

\subsection{投影对称群}

\begin{keypoints}
\item 投影对称群是对称性群的一个扩张。
\item 投影对称群对所有平均场相进行分类。
\item 只有当相应的平均场态稳定时，投影对称群才刻画物理自旋液体态中的量子序。
\end{keypoints}

让我们给出 PSG 的一个详细定义。PSG 是拟设的一种属性。它由保持拟设不变的所有变换构成。每个变换（即 PSG 中的每个元素）都可以写成一个对称性变换 $U$（如平移）与一个规范变换 $G_U$ 的组合。拟设在其 PSG 下的不变性可以表示为
\begin{align*}
G_U U(u_{\pmb{i}\pmb{j}}) &= u_{\pmb{i}\pmb{j}},\tag{9.4.1}\\
U(u_{\pmb{i}\pmb{j}}) &\equiv u_{U(\pmb{i}),U(\pmb{j})},\qquad G_U(u_{\pmb{i}\pmb{j}})\equiv G_U(\pmb{i})u_{\pmb{i}\pmb{j}}G_U^{\dagger}(\pmb{j}),\qquad G_U(\pmb{i})\in SU(2),
\end{align*}
对每个 $G_UU\in\text{PSG}$ 上式都成立。

每个 PSG 都包含一个特殊的子群，我们称之为不变规范群（invariant gauge group，IGG）。一个拟设的 IGG（记作 $\mathcal{G}$）由保持该拟设不变的所有纯规范变换构成：
\begin{equation*}
\mathcal{G}=\{W_{\pmb{i}}\,|\,W_{\pmb{i}}u_{\pmb{i}\pmb{j}}W_{\pmb{j}}^{\dagger}=u_{\pmb{i}\pmb{j}},\ W_{\pmb{i}}\in SU(2)\}
\tag{9.4.2}
\end{equation*}
如果我们想把 IGG 与一个对称性变换联系起来，那么相应的变换就是恒等对称性变换。

如果 IGG 是非平庸的，那么对一个固定的对称性变换 $U$，可以存在许多规范变换 $G_U$ 使得 $G_UU$ 保持拟设不变。若 $G_UU$ 在 $u_{\pmb{i}\pmb{j}}$ 的 PSG 中，则 $GG_UU$ 也在 PSG 中的充要条件是 $G\in\mathcal{G}$。这样，对每个对称性变换 $U$，$G_U$ 的不同选取与 IGG 中的元素一一对应。从上述定义我们看到，一个拟设的 PSG、IGG 与对称性群（symmetry group，SG）之间有如下关系：
\begin{equation*}
SG=PSG/IGG
\end{equation*}
这一关系告诉我们，PSG 是对称性群的一个投影表示，或称一个扩张（extension）。\footnote{更一般地，如果群 PSG 包含一个正规子群 IGG，使得 $PSG/IGG=SG$，我们就说群 PSG 是群 SG 的一个扩张。}

当然，两个规范等价拟设 $u_{\pmb{i}\pmb{j}}$ 与 $W(\pmb{i})u_{\pmb{i}\pmb{j}}W^{\dagger}(\pmb{j})$ 的 PSG 彼此相关。由 $WG_UU(u_{\pmb{i}\pmb{j}})=W(u_{\pmb{i}\pmb{j}})$，其中 $W(u_{\pmb{i}\pmb{j}})\equiv W(\pmb{i})u_{\pmb{i}\pmb{j}}W^{\dagger}(\pmb{j})$，我们发现
\begin{equation*}
WG_UUW^{-1}W(u_{\pmb{i}\pmb{j}})=WG_UW_U^{-1}UW(u_{\pmb{i}\pmb{j}})=W(u_{\pmb{i}\pmb{j}})
\end{equation*}
其中 $W_U\equiv UWU^{-1}$ 由 $W_U(\pmb{i})=W(U(\pmb{i}))$ 给出。这样，若 $G_UU$ 在拟设 $u_{\pmb{i}\pmb{j}}$ 的 PSG 中，则 $(WG_UW_U)U$ 在规范变换后的拟设 $W(\pmb{i})u_{\pmb{i}\pmb{j}}W^{\dagger}(\pmb{j})$ 的 PSG 中。我们看到，在规范变换 $W(\pmb{i})$ 之后，与对称性变换 $U$ 相联系的规范变换 $G_U$ 按如下方式改变：
\begin{equation*}
G_U(\pmb{i})\to W(\pmb{i})G_U(\pmb{i})W^{\dagger}(U(\pmb{i}))
\tag{9.4.3}
\end{equation*}

由于 PSG 是拟设的一种属性，我们可以把共享同一 PSG 的所有拟设归成一类。我们宣称，这样的类是一个普适类，对应于一个量子相。\footnote{更精确地说，这样的类由一个或若干个对应于量子相的普适类构成。（关于这一点的更详细讨论见 9.9 节。）}正是在这个意义上，我们说量子序由 PSG 刻画。

把量子序的 PSG 刻画与对称性破缺序的对称性刻画加以比较是很有启发的。我们知道，对称性破缺序可以由它的对称性性质描述。在数学上，我们说对称性破缺序由它的对称性群刻画。类似地，量子序也由群刻画。差别在于，量子序由 PSG——对称性群的扩张——刻画。我们看到，用投影对称群描述量子序，与用对称性群描述对称性破缺序，在概念上是相似的。我们还看到，具有相同对称性的量子态可以有许多不同的量子序，因为一个对称性群通常可以有许多不同的扩张。

除了能够刻画对称性破缺序之外，对称性破缺序的对称性描述还非常有用，因为它使我们无需了解系统的细节就能得到许多普适性质，例如无能隙 Nambu--Goldstone 模的数目。类似地，知道了一个量子序的 PSG，也使我们无需了解量子系统的细节就能得到它的低能性质。在 9.10 节中我们将证明，PSG 能够产生并保护低能规范涨落。事实上，低能规范涨落的规范群就是拟设的 IGG。这推广了 9.2.2 节得到的结果。除了无能隙规范玻色子之外，PSG 还能产生并保护无能隙费米子激发及其晶体动量。这些无能隙费米子甚至可以在纯玻色模型中产生。我们看到，无能隙规范玻色子与无能隙费米子可以有一个统一的起源——量子序。

这里我们想强调，PSG 真正分类的是不同的平均场相。连续的平均场相变总是伴随着 PSG 的改变。因此，严格说来，上述关于 PSG 及其物理含义的讨论只适用于平均场理论。然而，有些平均场态对涨落是稳定的。这些稳定平均场态的 PSG 就可以应用于相应的物理自旋液体，并描述这些自旋液体中真实的量子序。另一方面，也存在对涨落不稳定的平均场态。在这种情形下，相应的 PSG 可能不描述任何真实的量子序。

\subsection{对称 $Z_2$ 自旋液体的分类}

\begin{keypoints}
\item 我们分类了其 PSG 为格点对称性群（SG）之 $Z_2$ 扩张的所有平均场拟设，即 $PSG/Z_2=SG$。
\item 存在超过 100 种对称 $Z_2$ 自旋液体。
\item 不变 PSG 与代数 PSG 的概念。
\end{keypoints}

作为自旋液体量子序 PSG 理论的一个应用，我们想对与平移变换相联系的 PSG 进行分类。为简单起见，我们把 IGG $\mathcal{G}$ 限制为 $Z_2$。也就是说，我们想找出平移群的所有 $Z_2$ 扩张。

当 $\mathcal{G}=Z_2$ 时，它只含两个元素——规范变换 $G_1$ 与 $G_2$：
\begin{equation*}
\mathcal{G}=\{G_1,G_2\},\qquad G_1(\pmb{i})=\tau^0,\qquad G_2(\pmb{i})=-\tau^0.
\end{equation*}
一个拟设的 IGG 要为 $\mathcal{G}=Z_2$，该拟设就必须产生非共线的 $SU(2)$ 通量；否则 IGG 将大于 $Z_2$。拟设中的非共线 $SU(2)$ 通量把 $SU(2)$ 规范结构破缺到 $Z_2$ 规范结构，相应的自旋液体是一个 $Z_2$ 自旋液体。我们看到，平移群的 $Z_2$ 扩张在物理上对所有\emph{只}具有平移对称性的 $Z_2$ 自旋液体进行了分类。

考虑一个具有平移对称性的 $Z_2$ 自旋液体。作为平移群的 $Z_2$ 扩张，它的 PSG 由四个元素 $\pm G_xT_x$ 与 $\pm G_yT_y$ 生成\footnote{译注：原文作“$\pm G_yT_x$”，据式 (9.4.6) 及脚注（原书脚注 70）应为 $\pm G_yT_y$ 之误，此处按本意译出。}，其中
\begin{equation*}
T_x(u_{\pmb{ij}})=u_{\pmb{i}-\pmb{x},\,\pmb{j}-\pmb{x}},\qquad T_y(u_{\pmb{ij}})=u_{\pmb{i}-\pmb{y},\,\pmb{j}-\pmb{y}}.
\end{equation*}
由于拟设的平移对称性，我们可以选取一个规范，使拟设的所有 $SU(2)$ 通量算符都是平移不变的。也就是说，若两个回路 $C_1$ 与 $C_2$ 通过一个平移相联系，则 $P_{C_1}=P_{C_2}$。我们把这样的规范称为均匀规范（uniform gauge）。

在变换 $G_xT_x$ 下，以 $\pmb{i}$ 为基点的 $SU(2)$ 通量算符 $P_C$ 按如下方式变换：$P_C\to G_x(\pmb{i}')P_{T_x(C)}G_x^{\dagger}(\pmb{i}')\allowbreak=G_x(\pmb{i}')P_CG_x^{\dagger}(\pmb{i}')$，其中 $\pmb{i}'=T_x\pmb{i}$ 是平移后的回路 $T_x(C)$ 的基点。我们看到，$P_C$ 在均匀规范下的平移不变性要求：对任何回路 $C$ 都有 $G_x(\pmb{i}')P_CG_x^{\dagger}(\pmb{i}')=P_C$。同样，$G_y$ 也满足类似的条件。由于以同一基点为基点的不同 $SU(2)$ 通量算符互不对易，$G_{x,y}(\pmb{i})$ 在每个格点只能取 $\pm\tau^0$ 两个值之一。

我们注意到，形如 $W(\pmb{i})=\pm\tau^0$ 的依赖格点的规范变换不改变 $SU(2)$ 通量算符的平移不变性质。因此，我们可以利用这类规范变换，通过式 (9.4.3) 进一步简化 $G_{x,y}$。首先，我们可以选取规范使（见问题 9.4.1）
\begin{equation*}
G_y(\pmb{i})=\tau^0.
\tag{9.4.4}
\end{equation*}
我们注意到，满足 $W(\pmb{i})=W(i_x)$ 的规范变换不改变条件 $G_y(\pmb{i})=\tau^0$。我们可以利用这类规范变换使
\begin{equation*}
G_x(i_x,\ i_y=0)=\tau^0.
\tag{9.4.5}
\end{equation*}
由于 $x$ 与 $y$ 方向的平移相互对易，$G_{x,y}$ 必须满足（对任何拟设，无论是不是 $Z_2$）
\begin{equation*}
G_xT_xG_yT_y(G_xT_x)^{-1}(G_yT_y)^{-1}=G_xT_xG_yT_yT_x^{-1}G_x^{-1}T_y^{-1}G_y^{-1}\in\mathcal{G}.
\tag{9.4.6}
\end{equation*}
这意味着
\begin{equation*}
G_x(\pmb{i})G_y(\pmb{i}-\pmb{x})G_x^{-1}(\pmb{i}-\pmb{y})G_y(\pmb{i})^{-1}\in\mathcal{G}
\tag{9.4.7}
\end{equation*}
对 $Z_2$ 自旋液体，由于式 (9.4.4)，式 (9.4.7) 约化为
\begin{equation*}
G_x(\pmb{i})G_x^{-1}(\pmb{i}-\pmb{y})=+\tau^0
\tag{9.4.8}
\end{equation*}
或
\begin{equation*}
G_x(\pmb{i})G_x^{-1}(\pmb{i}-\pmb{y})=-\tau^0
\tag{9.4.9}
\end{equation*}
与式 (9.4.5) 和 (9.4.4) 联立，我们发现式 (9.4.8) 与 (9.4.9) 成为
\begin{equation*}
G_x(\pmb{i})=\tau^0,\qquad G_y(\pmb{i})=\tau^0
\tag{9.4.10}
\end{equation*}
以及
\begin{equation*}
G_x(\pmb{i})=(-)^{i_x}\tau^0,\qquad G_y(\pmb{i})=\tau^0
\tag{9.4.11}
\end{equation*}
这样，平移群只有两个规范不等价的 $Z_2$ 扩张。这两个 PSG 由 $G_xT_x$ 与 $G_yT_y$ 生成，其中的 $G_{x,y}$ 由式 (9.4.10) 与 (9.4.11) 给出。\footnote{我们想指出，若 $G_x$ 与 $G_y$ 由式 (9.4.8) 给出，则 $G_xT_x$ 与 $G_yT_y$ 满足通常的平移代数 $G_xT_xG_yT_y=G_yT_yG_xT_x$。当 $G_x$ 与 $G_y$ 由式 (9.4.9) 给出时，$G_xT_x$ 与 $G_yT_y$ 满足磁平移代数 $G_xT_xG_yT_y=-G_yT_yG_xT_x$。}在 PSG 分类下，\emph{只}具有平移对称性而没有其他对称性的 $Z_2$ 自旋液体只有两种类型。在 PSG (9.4.10) 下不变的拟设具有形式
\begin{equation*}
u_{\pmb{i},\pmb{i}+\pmb{m}}=u_{\pmb{m}}
\tag{9.4.12}
\end{equation*}
而在 PSG (9.4.11) 下不变的拟设具有形式
\begin{equation*}
u_{\pmb{i},\pmb{i}+\pmb{m}}=(-)^{m i_x}u_{\pmb{m}}
\tag{9.4.13}
\end{equation*}

通过上面的例子我们看到，PSG 是一个非常强有力的工具。对于具有给定对称性和低能规范结构的（平均场）自旋液体，它可以给出完全的分类。

在上文中，我们研究了只具有平移对称性而无其他对称性的 $Z_2$ 自旋液体，发现这样的自旋液体只有两种类型。然而，如果自旋液体具有更多对称性，类型就会多得多。Wen (2002c) 利用 PSG 得到了对称 $Z_2$ 自旋液体的分类。这一分类来自如下观察：规范变换 $G_{x,y}$、$G_{P_x,P_y,P_{xy}}$ 与 $G_T$ 必须满足某些与式 (9.4.7) 类似的代数关系。求解这些代数关系并剔除规范等价的解，我们发现，由平移、宇称与时间反演变换 $\{T_{x,y},P_{x,y,xy},T\}$ 生成的对称性群共有 272 个不同的 $Z_2$ 扩张。这里我们仅列出这 272 个 $Z_2$ PSG。这些 PSG 由 $(G_xT_x,\ G_yT_y,\ G_TT,\ G_{P_x}P_x,\ G_{P_y}P_y,\ G_{P_{xy}}P_{xy})$ 生成。PSG 可以分成两类。第一类为
\begin{equation*}
\begin{array}{rcl}
G_x(\pmb{i})=\tau^0, & & G_y(\pmb{i})=\tau^0\\[2pt]
G_{P_x}(\pmb{i})=\eta_{xpx}^{i_x}\eta_{xpy}^{i_y}g_{P_x}, & & G_{P_y}(\pmb{i})=\eta_{xpy}^{i_x}\eta_{xpx}^{i_y}g_{P_y}\\[2pt]
G_{P_{xy}}(\pmb{i})=g_{P_{xy}}, & & G_T(\pmb{i})=\eta_t^{\pmb{i}}g_T
\end{array}
\tag{9.4.14}
\end{equation*}
第二类为
\begin{equation*}
\begin{array}{rcl}
G_x(\pmb{i})=(-)^{i_x}\tau^0, & & G_y(\pmb{i})=\tau^0\\[2pt]
G_{P_x}(\pmb{i})=\eta_{xpx}^{i_x}\eta_{xpy}^{i_y}g_{P_x}, & & G_{P_y}(\pmb{i})=\eta_{xpy}^{i_x}\eta_{xpx}^{i_y}g_{P_y}\\[2pt]
G_{P_{xy}}(\pmb{i})=(-)^{i_x i_y}g_{P_{xy}}, & & G_T(\pmb{i})=\eta_t^{\pmb{i}}g_T
\end{array}
\tag{9.4.15}
\end{equation*}
这里三个 $\eta$ 可以各自独立地取 $\pm1$ 两个值。$g$ 有 17 种不同的选择，列于表 9.1 之中。因此，共有 $2\times17\times2^{3}=272$ 个不同的 PSG。它们可能给出二维正方格点上 272 种不同类型的对称 $Z_2$ 自旋液体。

为了标记这 272 个 PSG，我们提出如下方案：
\begin{equation*}
\text{Z2A}(g_{px})_{\eta_{xpx}}(g_{py})_{\eta_{xpy}}g_{pxy}(g_t)_{\eta_t},
\tag{9.4.16}
\end{equation*}
\begin{equation*}
\text{Z2B}(g_{px})_{\eta_{xpx}}(g_{py})_{\eta_{xpy}}g_{pxy}(g_t)_{\eta_t}.
\tag{9.4.17}
\end{equation*}
标记 Z2A… 对应式 (9.4.14) 的情形，标记 Z2B… 对应式 (9.4.15) 的情形。一个典型的标记形如 $\text{Z2A}\tau_{+}^{1}\tau_{-}^{2}\tau^{12}\tau_{-}^{3}$。我们还将使用一种缩写记号：把 $(\tau^{0},\tau^{1},\tau^{2},\tau^{3})$ 或 $(\tau_{+}^{0},\tau_{+}^{1},\tau_{+}^{2},\tau_{+}^{3})$ 换成 $(0,1,2,3)$，把 $(\tau_{-}^{0},\tau_{-}^{1},\tau_{-}^{2},\tau_{-}^{3})$ 换成 $(n,x,y,z)$。例如，$\text{Z2A}\tau_{+}^{1}\tau_{-}^{0}\tau^{12}\tau_{-}^{3}$ 可缩写为 Z2A1n(12)z。

严格说来，这 272 个不同的 $Z_2$ PSG 是所谓的代数 PSG（algebraic PSG）。代数 PSG 定义为对称性群的扩张。代数 PSG 不同于不变 PSG（invariant PSG），后者定义为保持一个拟设 $u_{\pmb{ij}}$ 不变的所有变换的集合。虽然不变 PSG 必定是代数 PSG，但代数 PSG 未必是不变 PSG。这是因为某些代数 PSG 具有如下性质：在代数 PSG $G_1$ 下不变的\emph{任何}拟设 $u_{\pmb{ij}}$，实际上可能在一个更大的 PSG $G_2$ 下也不变。这时，原来的 PSG $G_1$ 不可能是该拟设的不变 PSG，拟设的不变 PSG 实际上由更大的 PSG $G_2$ 给出。如果我们把自己限制在通过拟设 $u_{\pmb{ij}}$ 构造的自旋液体上，那么就应当舍弃那些不是不变 PSG 的代数 PSG，因为这些代数 PSG 并不刻画平均场自旋液体。

\begin{center}
\small 表 9.1\quad $Z_2$ PSG 的 $g_{P_{xy}}$、$g_{P_x}$、$g_{P_y}$ 与 $g_T$ 的 17 种选择。这里 $\tau^{ab}=(\tau^{a}+\tau^{b})/\sqrt{2}$，$\tau^{a\bar{b}}=(\tau^{a}-\tau^{b})/\sqrt{2}$。\\[6pt]
\begin{tabular}{cccccccc}
\toprule
$g_{P_{xy}}$ & $g_{P_x}$ & $g_{P_y}$ & $g_T$ & $g_{P_{xy}}$ & $g_{P_x}$ & $g_{P_y}$ & $g_T$ \\
\midrule
$\tau^0$ & $\tau^0$ & $\tau^0$ & $\tau^0$ & $\tau^0$ & $\mathrm{i}\tau^3$ & $\mathrm{i}\tau^3$ & $\tau^0$ \\
$\mathrm{i}\tau^3$ & $\tau^0$ & $\tau^0$ & $\tau^0$ & $\mathrm{i}\tau^3$ & $\mathrm{i}\tau^3$ & $\mathrm{i}\tau^3$ & $\tau^0$ \\
$\mathrm{i}\tau^3$ & $\mathrm{i}\tau^1$ & $\mathrm{i}\tau^1$ & $\tau^0$ & $\tau^0$ & $\tau^0$ & $\tau^0$ & $\mathrm{i}\tau^3$ \\
$\tau^0$ & $\mathrm{i}\tau^3$ & $\mathrm{i}\tau^3$ & $\mathrm{i}\tau^3$ & $\tau^0$ & $\mathrm{i}\tau^1$ & $\mathrm{i}\tau^1$ & $\mathrm{i}\tau^3$ \\
$\mathrm{i}\tau^3$ & $\tau^0$ & $\tau^0$ & $\mathrm{i}\tau^3$ & $\mathrm{i}\tau^3$ & $\mathrm{i}\tau^3$ & $\mathrm{i}\tau^3$ & $\mathrm{i}\tau^3$ \\
$\mathrm{i}\tau^3$ & $\mathrm{i}\tau^1$ & $\mathrm{i}\tau^1$ & $\mathrm{i}\tau^3$ & $\mathrm{i}\tau^1$ & $\tau^0$ & $\tau^0$ & $\mathrm{i}\tau^3$ \\
$\mathrm{i}\tau^1$ & $\mathrm{i}\tau^3$ & $\mathrm{i}\tau^3$ & $\mathrm{i}\tau^3$ & $\mathrm{i}\tau^1$ & $\mathrm{i}\tau^1$ & $\mathrm{i}\tau^1$ & $\mathrm{i}\tau^3$ \\
$\mathrm{i}\tau^1$ & $\mathrm{i}\tau^2$ & $\mathrm{i}\tau^2$ & $\mathrm{i}\tau^3$ & $\mathrm{i}\tau^{12}$ & $\mathrm{i}\tau^1$ & $\mathrm{i}\tau^2$ & $\mathrm{i}\tau^0$ \\
$\mathrm{i}\tau^{12}$ & $\mathrm{i}\tau^1$ & $\mathrm{i}\tau^2$ & $\mathrm{i}\tau^3$ & & & & \\
\bottomrule
\end{tabular}
\end{center}

我们发现，在 272 个代数 $Z_2$ PSG 中，至少有 76 个不是不变 PSG。因此，272 个代数 $Z_2$ PSG 至多能给出 196 种可能的 $Z_2$ 自旋液体。

\subsection{一个 $Z_2$ 与一个 $U(1)$ 投影对称群及其拟设}

\begin{keypoints}
\item 对给定的 PSG，我们可以找出其最普遍的拟设。这些拟设共有的物理性质，就是与该 PSG 相联系的普适性质。
\end{keypoints}

作为 $Z_2$ PSG 的一个应用，让我们构造在 Z2A0013 PSG 下不变的最普遍拟设：
\begin{equation*}
\begin{array}{rcl}
G_x(\pmb{i})=\tau^0, & & G_y(\pmb{i})=\tau^0\\[2pt]
G_{P_x}(\pmb{i})=\tau^0, & & G_{P_y}(\pmb{i})=\tau^0\\[2pt]
G_{P_{xy}}(\pmb{i})=\mathrm{i}\tau^1, & & G_T(\pmb{i})=\mathrm{i}\tau^3
\end{array}
\tag{9.4.18}
\end{equation*}
由 $G_xT_x$ 与 $G_yT_y$ 下的不变性，再考虑到 $G_x$ 与 $G_y$ 是平庸的，Z2A0013 拟设直接就是平移不变的，形式为 $u_{\pmb{i},\pmb{i}+\pmb{m}}=u_{\pmb{m}}^{\mu}\tau^{\mu}$。求出最普遍 Z2A0013 拟设的关键是

% ---- 块界备注：上句延续至书页 406（p421）首行 "find out which $u_m^\mu$ must vanish."
% ---- （"……找出哪些 $u_{\pmb{m}}^{\mu}$ 必须为零。"），请 ch9_e 以 %JOIN% 续接，勿重复勿遗漏。

% ---------------- 新增术语（ch9_d） ----------------
% | extension (of a group) | 扩张 |
% | projective representation | 投影表示 |
% | normal subgroup | 正规子群 |
% | invariant group | 不变群 |
% | symmetry group (SG) | 对称性群（SG） |
% | invariant PSG | 不变 PSG |
% | algebraic PSG | 代数 PSG |
% | uniform gauge | 均匀规范 |
% | magnetic translational algebra | 磁平移代数（ch7_a 已收“磁平移”） |
% | sector | 扇区 |
% | pure gauge transformation | 纯规范变换 |
% | collinearness | 共线性 |
% | thermodynamic limit | 热力学极限 |
% | flux conservation | 通量守恒 |
% | infra-red divergence | 红外发散 |
% | quantum phase | 量子相 |
% | universal property | 普适性质 |
% | topologically-degenerate ground state | 拓扑简并基态 |
% | first-order mean-field theory | 一阶平均场理论 |
% | Z2-gapped / Z2-gapless / Z2-linear spin liquid | Z2 有能隙 / 无能隙 / 线性自旋液体 |
% | SU(2)-gapless / SU(2)-linear state | SU(2) 无能隙态 / SU(2) 线性态 |
% | U(1)-linear state | U(1) 线性态 |
% | invariant gauge group (IGG) | 不变规范群（IGG）（术语表已收，本块首次正文出现） |
% | compactified lattice | 紧致化格点（compactification 紧致化已收） |
